\chapter{Launcher App}

\section{Problem}
Es hat sich im Laufe der Projektentwicklung bei mehreren AR-Projekten herausgestellt, dass die geplante Idee mit einer Webanwendung nicht möglich ist. Durch diese Gegebenheit sind unterschiedliche Ideen für die Bereitstellung für die Besucher*innen entstanden. Eine dieser Ideen war, dass vorne Tablets für die Besucher*innen herausgegeben werden können und diese Idee ist im folgenden weiter verfolgt worden.

\section{Idee}
Diese Idee besteht daraus, dass Tablets vorne an der Rezeption liegen, auf denen die hier entwickelten Projekte und auch zukünftige Projekte drauf installiert sind. Diese können die Besucher*innen vorne ausleihen und damit die XR interaktiven Elemente des Museums erleben. Wenn das Tablet angeschaltet wird, öffnet sich standartmäßig eine Launcher-App, welche für jedes Projekt einen Button anzeigt, um diese zu öffnen. Eine Bedingung ist, dass jede App einen Weg liefern muss, zurück zur Launcher-App zu kommen.

\section{Projekt}
Im Rahmen dieser Idee, ist ein Proof of concept mit der OCMExplore-Gruppe entstanden. Dieses bietet einen Button in der Mitte des Bildschirms, mit welchem die OCMExplore-App geöffnet werden kann.

\section{Aufgabeneinteilung}
Dieses Projekt ist von Louisa Valerie Müller entwickelt worden.

\section{Weitere Überlegungen}
Im laufe der Entwicklung des Proof of concepts sind weitere Überlegungen entstanden.

\subsection{Modularität}
Dadurch, dass die technische Grundlage rein programmativer Natur ist, bietet sich die Möglichkeit ein Adminpanel einzubinden, in welchem weitere Projekte eingebunden werden können. Dies würde die programmative Wartungsarbeit stark reduzieren.

\subsection{Vor-Ort Umsetzung}
Die Umsetzung vor Ort bestünde aus Android-Tablets, welche XR-Kompatibilität bieten. Dazu kann man dann spezialisiert auf diese Tablets entwickeln, welche genauere Vorraussetzungen dann noch erstellen würde. Diese Tablets sind so eingerichtet, dass die Besucher:innen nur auf diese Launcher-App zugreifen können. Dies könnte mit eine App, wie zum Beispiel, MacroDroid umgesetzt werden. Vorne an der Rezeption werden Tablets verliehen und auch geladen. Genauere Konfiguration von Android muss weiter ausprobiert werden. Für die Zuordnung von App und Ort im Museum könnte ein Iconsystem dienen, wo für jede App ein Icon existiert und diese dann bei den jeweiligen Stationen platziert sind.

\subsection{Design}
Das Design der App richtet sich nach dem Stil des OCM's. Genaueres steht im technischen Bericht.

\section{Fazit}
Die Idee einer allgemeinen Launcher App ist technisch umsetzbar und von diesem Proof of concept aus weiter entwickelbar.